\lecture{11}{2026-10-02}{}

\section*{Principal branches}

How to define \(\sqrt{z}\)? It's the inverse of the squaring map \(z \mapsto z^2\), but this is not an injective map in \(\mathbb{C}\).

Given \(z \in \mathbb{C} \setminus \mathbb{R}_{\le 0}\), we can write \(z = r e^{i \theta}\) with \(r > 0\) and \(\theta \in (-\pi, \pi)\).
The \emph{principal argument} of \(z\) is \(\operatorname{Arg}(z) = \theta\).

The \emph{principal branch square-root} is the analytic function
\begin{align*}
    \mathbb{C} \setminus \mathbb{R}_{\le 0} &\longrightarrow \mathbb{C} \\
    z &\mapsto \sqrt{|z|} e^{i \operatorname{Arg}(z)/2}.
\end{align*}

Note we cannot extend this definition to the entire \(\mathbb{C}\), since
\[
\lim_{\theta \to \pi} \sqrt{r e^{i \theta}} = \sqrt{r} e^{i \pi/2} = i \sqrt{r} \neq -i \sqrt{r} = \lim_{\theta \to -\pi} \sqrt{r e^{i \theta}}.
\]

Removing \(\mathbb{R}_{\le 0}\) from the domain of \(\sqrt{z}\) is an example of a \emph{branch cut}.

We can remove any ray starting from the origin from \(\mathbb{C}\) to get a well-defined branch of the square-root function.
For example, if we remove the negative imaginary ray \(\mathbb{R}_{\le 0} i\), we can define a branch of the square-root function on \(\mathbb{C} \setminus \mathbb{R}_{\le 0} i\) by
\[
z = r e^{i \theta} \mapsto \sqrt{|z|} e^{i \theta/2}, \quad \theta \in (-\pi/2, 3\pi/2).
\]

Defining \(z^r\) for \(r \in \mathbb{Q}\) is similar.

We can also define the principal branch of the logarithm function \(\log(z)\) on \(\mathbb{C} \setminus \mathbb{R}_{\le 0}\) by
\[
\log(z) = \log(|z|) + i \operatorname{Arg}(z).
\]

A \emph{slit disk} is a disk with a ray from the center removed.
The function \(f(z)\) has an \emph{algebraic singularity} (or \emph{algebraic branch point}) at \(z = a \in \mathbb{C}\) if
\[
    f(z) = \sum_{j \in N} c_j \left( (z-a)^{1/R} \right)^j
\]
in a slit disk around \(z = a\) of radius \(R\), for some \(R \in \mathbb{Z}_{>1}\), finite \(N \subseteq \mathbb{Z}\), and \(R > 0\), where the map \(t \mapsto t^{1/R}\) a branch of the \(R\)-th root function on a slit disk around \(t = 0\) of radius \(R\), and there exists \(c_n \neq 0\) where \(n / R \notin \mathbb{Z}\).
% what are these conditions?
Equivalently, \(f(z)\) is not analytic at \(z = a\), but \(f(z^R + a)\) is analytic at \(z = 0\) for some integer \(R > 1\).

Similarly, \(f(z)\) has a \emph{logarithmic singularity} (or \emph{logarithmic branch point}) at \(z = a\) if
\[
    f(z) = (z-a)^N \sum_{j=0}^d c_j \left( \log(z-a) \right)^j
\]
in a slit disk around \(z = a\) of radius \(R\), for some \(N \in \mathbb{Z}\), \(d \in \mathbb{Z}_{\ge 0}\), and \(R > 0\), where the map \(t \mapsto \log(t)\) is a branch of the logarithm function on a slit disk around \(t = 0\) of radius \(R\), and there exists \(c_d \neq 0\).

Given a branch of the logarithm function \(\log(z)\), we can define
\[
    z^\alpha = e^{\alpha \log(z)}
\]
for any \(\alpha \in \mathbb{C}\).
If \(\alpha \in \mathbb{Z}_{\ge 0}\), then \(z^\alpha\) is has no singularity at \(z = 0\).
If \(\alpha \in \mathbb{Z}_{< 0}\), then \(z^\alpha\) has a pole at \(z = 0\).

\section*{Transfer theorems}

If \(\alpha \in \mathbb{Z}_{\ge 0}\), then
\[
    [z^n] (1 - z)^{-\alpha} = \binom{n + \alpha - 1}{n} = \frac{n^{\alpha - 1}}{(\alpha-1)!} \left( 1 + O\left( \frac{1}{n} \right) \right).
\]

We can generalize this to any \(\alpha \in \mathbb{C} \setminus \mathbb{Z}_{\le 0}\) using the Gamma function.

\begin{definition}[Gamma function]
    The \emph{Gamma function} is defined for \(\alpha \in \mathbb{C} \setminus \mathbb{Z}_{\le 0}\) by
    \[
        \Gamma(\alpha) = \int_0^\infty t^{\alpha - 1} e^{-t} dt.
    \]
\end{definition}

It has simple poles at \(\alpha \in \mathbb{Z}_{\le 0}\) and satisfies the functional equation \(\Gamma(\alpha + 1) = \alpha \Gamma(\alpha)\).

A useful value is \(\Gamma(1/2) = \sqrt{\pi}\).

\begin{theorem}[Transfer theorems]
    \begin{itemize}
        \item \(f(z) \sim a(z)\) as \(z \to 1\) implies \([z^n] f(z) \sim [z^n] a(z)\) as \(n \to \infty\).
        \item \(f(z) = a(z) + O(b(z))\) as \(z \to 1\) implies \([z^n] f(z) = [z^n] a(z) + O([z^n] b(z))\) as \(n \to \infty\).
    \end{itemize}
\end{theorem}

\begin{lemma}
    \(\alpha \in \mathbb{C} \setminus \mathbb{Z}_{\le 0}\) and \(F(z) = (1 - z)^{-\alpha}\) implies \([z^n] F(z) \sim \frac{n^{\alpha - 1}}{\Gamma(\alpha)}\) as \(n \to \infty\).
\end{lemma}